% Working supplement. Does not replace registro_k.tex or settle E108 by fiat.
% Provenance and scope: INFORME_RECUPERACION.md, same folder.
\subsection{Incidence evaluation preceding the dodecaphasic register}
\label{subsec:registro-evaluacion-incidencial}

Recovering a register from its cyclic differences and
total charge is a different question from evaluating those coordinates
on the incidence register. The first problem is solved by the
transverse extractor of Proposition~\ref{prop:k-transversal}.
For the second, a historical presentation of the
register provides three integer counts in each window. It is useful
to make their composition explicit and determine exactly the information
required for their evaluation.

Let \(\mathscr L\) be a previously generated ledger of visits and oriented
traces. Each visit retains the path, APP position, sheet,
time, multiplicity, and boundary incidence. Arithmetic
evaluation produces the integer and its decomposition
\[
 n_t=r_t+9q_t,\qquad 1\leq r_t\leq9.
\]
The integer \(n_t\) is computed from the sum or product sheet; its
residue is not received as an independent value. The distinction between
corona incidence and interior incidence remains explicit for
visits with residue \(9\).

In window \(E_m\), the count-based presentation considers
\begin{align}
 A_m&=\sum_{v\in E_m}w(v)
             \mathbf1_{\{1,3,5,7\}}(r(v)),\\
 C_m&=\sum_{j\geq0}\bigl(\nu^-_{120,m}(j)-\nu^+_{120,m}(j)\bigr),\\
 V_m&=\sum_{v\in E_m}w(v)\mathbf1_{\{9\}}(r(v))\epsilon_\partial(v).
\end{align}
Here \(w(v)\) is the incidence multiplicity, and
\(\epsilon_\partial(v)=1\) on the corona, \(-1\) in the interior.
The traces counted by \(\nu^\pm\) must be enumerated by the
corresponding rule in the article. The integer sum requires finite support or an
alternative construction determining its meaning. The number of
chronological incidences and the reduced length \(120(1+3j)\) have
different types; relating them requires the reader's unit of length.

Define \([x]_{10}\in\{0,\ldots,9\}\) as the Euclidean
representative, retaining the quotients
\[
 x=[x]_{10}+10\lfloor x/10\rfloor.
\]
The incidence block is
\begin{equation}
 z_m=100[A_m]_{10}+10[C_m]_{10}+[V_m]_{10}.
 \label{eq:registro-bloque-incidencial}
\end{equation}
Composition yields the map
\begin{equation}
 \mathcal E_{\rm cnt}(\mathscr L)=
 \left((z_m-z_{m+3})_{m=1}^{12},
       (z_m-z_{m+4})_{m=1}^{12},\sum_{m=1}^{12}z_m\right).
\label{eq:registro-composicion-incidencial}
\end{equation}
Indices are read cyclically modulo \(12\). In this subsection,
\((D_jz)_m=z_m-z_{m+j}\), for \(j=3,4\), and
\(Q(z)=\sum_{m=1}^{12}z_m\); hence the preceding composition is
\((D_3z,D_4z,Q(z))\). This convention is retained in
\S\ref{subsec:k-transversal}, where complete integral
reconstruction is proved.
This formula produces its coordinates from the counts. It does not use
\(K\), \(U\), \(\alpha\), or the expected transverse coordinates
as arguments. Identifying it with
\(\mathcal E_{108}^{90,120}\) requires composing it with the actual collection and
incidences of the terminal state; this does not follow solely
from invertibility of the transverse system.

\begin{proposition}[Necessary and sufficient arithmetic information]
\label{prop:registro-36-residuos}
Let \(N,N'\in\mathbb Z^{12\times3}\) be ordered by the counts
\((A,C,V)\). For the map defined by
\eqref{eq:registro-bloque-incidencial}--\eqref{eq:registro-composicion-incidencial},
\[
 \mathcal E_{\rm cnt}(N)=\mathcal E_{\rm cnt}(N')
 \quad\Longleftrightarrow\quad
 N-N'\in10\mathbb Z^{36}.
\]
In particular, the \(36\) sectorial residues determine exactly
the output of \(25\) coordinates. Retention of the quotients and
provenance constitutes additional information in the register.
\end{proposition}

\begin{proof}
Equality modulo \(10\) produces the same blocks and hence
the same transverse coordinates. Conversely, if the transverse
coordinates agree, \(D_3(z-z')=D_4(z-z')=0\). The steps
\(3\) and \(4\) generate the cycle of length \(12\), so
\(z-z'\) is constant. Equality of total charges makes that
constant zero. Finally, the decimal representation of an integer between
\(0\) and \(999\) has three unique digits. This recovers the
three residues of each window. The statement is an algebraic identity
for every \(N,N'\), not an inference drawn from finite tests.
\end{proof}

\subsection{Conjugation of counts and boundary terms}
\label{subsec:registro-conjugacion-incidencial}

Let \(\rho(m)=13-m\). Consider an incidence conjugation
that reverses the order of the windows, preserves the arithmetic and
boundary classes of visits, and exchanges the trace channels.
Then
\[
 (A_m^\dagger,C_m^\dagger,V_m^\dagger)
 =(A_{\rho(m)},-C_{\rho(m)},V_{\rho(m)}).
\]
If \(c_m=[C_m]_{10}\), decimal reduction produces the identity
\begin{equation}
 z_m^\dagger=z_{\rho(m)}
       +100\mathbf1_{\{c_{\rho(m)}\ne0\}}-20c_{\rho(m)}.
 \label{eq:registro-conjugacion-decimal}
\end{equation}
Indeed, \([-C]_{10}=0\) when \([C]_{10}=0\), and
\([-C]_{10}=10-[C]_{10}\) otherwise. The formula retains
exactly that difference. Negating the channel is not equivalent to negating
the entire decimal block.

Applying this identity to a concrete TPK involution requires
checking its action on incidences. For an occupancy reader
evaluated before the advance, reversal of a path
\(\gamma=(x_0,\ldots,x_n)\) satisfies
\[
 \sum_{k=0}^{n-1}f(x_{n-k})
 =\sum_{k=0}^{n-1}f(x_k)+f(x_n)-f(x_0).
\]
The endpoint terms are incorporated before reduction modulo
\(10\). They vanish on closed paths; on open windows they may
alter the residues. This correction, already present in the bidirectional
development of transport, specifies what the connection
between the path and the counts must preserve.

\begin{proposition}[Necessity of the register's nonperiodic information]
Suppose that the observable of a fixed collection of paths satisfies
\(o(t+54)=o(t)\), that multiplicity is preserved under this return,
and that the same local reader acts on the windows of \(9\) instants.
Then its counts and blocks satisfy
\[
 (A,C,V)_{m+6}=(A,C,V)_m,\qquad z_{m+6}=z_m.
\]
\end{proposition}

\begin{proof}
The translation \(t\mapsto t+54\) maps \(E_m\) bijectively
onto \(E_{m+6}\), preserving each counted incidence and its multiplicity.
The equality passes to the sum and to decimal reduction.
\end{proof}

Thus a presentation with \(12\) distinct blocks requires the
reader to retain some information that does not factor through that observable
of period \(54\): incidence selection, effective orientation,
endpoints, multiplicity, or memory. The proposition identifies an inadmissible
reduction of the state; it does not assert that the complete dynamics has
period \(54\), nor does it determine its terminal reader by itself.
